\documentclass[a4paper,11pt]{article}
\usepackage{amsmath,amssymb,amsthm, tikz,titlesec,hyperref,esint,braket, graphicx, mathrsfs, enumitem}
\usepackage[a4paper,margin=2cm]{geometry}
\linespread{1.3}
\newtheorem{claim}{Claim}[section]

\newcommand\name{Park Chanwoo}   % Name of the student
\newcommand\university{KAIST} % Name of the university
\newcommand\department{Physics} % Name of the department
\newcommand\studentid{20230297} % Student ID
\newcommand\s{\,\;}


\newenvironment{solution}[1]
  {\renewcommand\qedsymbol{$\square$}\begin{proof}[\textbf{Solution#1}]}
  {\end{proof}}
\newenvironment{note}
  {\renewcommand\qedsymbol{$\blacksquare$}\begin{proof}[\textnormal{\textbf{note}}]}
  {\end{proof}}

\title{KAIST\\2026 Special Topics in Physics Quantum Electronic Devices I: Interferometers and Fractional Particles\\
Homework 1\bigskip}
\author{\textbf{\Large \name} \\
% University: \university\\
Department: \department\\
Student ID: \studentid}
\date{\today}

\begin{document}
\thispagestyle{empty}
\maketitle
\tableofcontents
\titleformat{\section}[frame]{\pagebreak}{\filright
\footnotesize  \enspace \textsf{KAIST --- Interferometers and Fractional Particles 2026 Spring}\enspace}{6pt}{\Large\bfseries\filcenter}

\newcommand{\tr}{\operatorname{tr}}
\newcommand{\sgn}{\operatorname{sgn}}
\newcommand{\supp}{\operatorname{supp}}
\renewcommand{\Re}{\operatorname{Re}}
\renewcommand{\Im}{\operatorname{Im}}
\newcommand{\boltz}{k_{\mathrm{B}}}

\section{\#1}

Assume $n(t)$ obeys the diffusion equation and the initial condition $n(x,0) = \delta(x)$, Then we can solve the diffusion equation by Fourier transform as follows:
\begin{gather}
    \partial_t n(x, t) = D \nabla^2 n(x, t) \implies \partial_t \tilde{n}(k, t) = -D k^2 \tilde{n}(k, t) \\
    \tilde{n}(k, t) = \tilde{n}(k, 0) e^{-D k^2 t} = e^{-D k^2 t} \\
    n(x,t) = \frac{1}{(2\pi)^d}\int e^{ik\cdot x} e^{-Dk^2 t} dk= \frac{1}{(4\pi Dt)^{d/2}} \exp\left(-\frac{x^2}{4Dt}\right)
\end{gather}
Due to the conservation law of diffusion equation, the denominator is 1:
\begin{equation}
    \int n(x,t) dx = 1
\end{equation}
And the mean absolute displacement can be evaluated as follows:
\begin{equation}
    \langle |x| \rangle = \int |x| n(x,t) dx = \frac{1}{(4\pi Dt)^{d/2}} \int_0^\infty r \exp\left(-\frac{r^2}{4Dt}\right) \Omega_d r^{d-1} dr = \frac{2 \Gamma(\frac{d+1}{2})}{\Gamma(\frac{d}{2})} \sqrt{Dt}
\end{equation}
Where $\Omega_d$ is the surface area of the unit sphere in $d$-dimension.

Hence,
\begin{equation}
    l(t) \propto \sqrt{Dt}
\end{equation}

\section{\#2}

The total number of particle can be evaluated by integrating the density of states as follows:
\begin{equation}
    N = \sum_{\frac{2\pi}{L}\sqrt{n_x^2 + n_y^2} \le k_F} \sum_{\sigma} 1 \approx \int_{k_x^2 + k_y^2 <= k_F^2} g_s\left(\frac{L}{2\pi}\right)^2 dk_x dk_y = \frac{L^2 g_s k_F^2}{4\pi}
\end{equation}
Using the relation $L\times L = V$ and $n_e = N/V$, 
\begin{eqnarray}
    n_e = \frac{g_s k_F^2}{4\pi}
\end{eqnarray}


\section{\#3}

Assume that the dispersion relation of the system is given by $E(k) = f(|k|)$. Then the density of states per area can be derived as follows:
\begin{gather}
    \sum_{\sigma} \left(\frac{L}{2\pi}\right)^2 dk_x dk_y = 2g_s\pi k \left(\frac{L}{2\pi}\right)^2dk = \frac{2g_s\pi k}{f'(k)} \left(\frac{L}{2\pi}\right)^2 dE = A\rho(E) dE \\
    \implies \rho(E) = \frac{g_s k}{2\pi f'(k)}
\end{gather}


For the quadratic dispersion relation, $f(k) = \frac{\hbar^2 k^2}{2m}$,
\begin{equation}
    \rho(E) = \frac{g_s k}{2\pi f'(k)} = \frac{g_s m}{2\pi \hbar^2}
\end{equation}


For the linear dispersion relation, $f(k) = \hbar v k$,
\begin{equation}
    \rho(E) = \frac{g_s k}{2\pi f'(k)} = \frac{g_s k }{2\pi \hbar v} = \frac{g_s E}{2\pi (\hbar v)^2}
\end{equation}


\end{document}
